\chapter{Command-Line Reference}
\label{ch:commands}

The exact parser generated by the installed binary is authoritative. Use
\command{blackship --help} and \command{blackship <command> --help} when an
option differs from this edition.

\section{Global options}

\begin{shiptable}{L{0.30\textwidth}L{0.60\textwidth}}
\toprule
Option & Effect \\
\midrule
\command{-c, --config <path>} & Select the manifest; default is
\file{blackship.toml}. \\
\command{-v, --verbose} & Emit additional operational diagnostics. \\
\command{--no-escalate} & Fail privileged commands rather than re-executing
through mdo, doas, or sudo. \\
\command{--non-interactive} & Add non-prompting behavior to doas/sudo; mdo is
inherently policy-based and non-interactive. \\
\command{--format text|json} & Select human text or the versioned machine
contract for supported commands. \\
\bottomrule
\end{shiptable}

Global options precede the subcommand. Use \command{--} to terminate option
parsing before an in-jail argv.

\section{Lifecycle and host commands}

\begin{shiptable}{L{0.39\textwidth}L{0.51\textwidth}}
\toprule
Command & Purpose \\
\midrule
\command{up [target] [--all] [--dry-run]} & Start selected jails with their
dependencies. \\
\command{down [target] [--all] [--dry-run]} & Stop in cascade/reverse dependency
order. \\
\command{restart [target] [--all] [--dry-run]} & Compose stop and start plans. \\
\command{eva [target] [--all] [--dry-run]} & Apply live-safe manifest drift. \\
\command{autostart [--dry-run]} & Reconcile running project jails to the
reserved \setting{autostart} tag and its dependency closure. \\
\command{ps [--host] [--json]} & List configured jail state and JID. With
\command{--host}, also include running jails outside the selected manifest. \\
\command{check} & Validate and resolve configuration; warn on releases not bootstrapped. \\
\command{capabilities} & Report binary, kernel, feature, facility, and escalator
capabilities. \\
\command{host} & Report the comprehensive host planning profile as JSON. \\
\command{setup [--enable-pf]} & Recreate networks/anchors and prepare the host;
optionally configure and start PF. \\
\command{cleanup <jail> [--force]} & Remove resources after a failed lifecycle. \\
\command{gc [jail] [--dry-run] [--force] [--prune-unreferenced]} & Find and
reclaim orphaned resources, including datasets left by a restarted jail and
empty bridges Blackship created; pruning broadens ownership beyond tracked
state. \\
\bottomrule
\end{shiptable}

\section{Execution and files}

\begin{shiptable}{L{0.39\textwidth}L{0.51\textwidth}}
\toprule
Command & Purpose \\
\midrule
\command{console <jail> [--user user]} & Open an interactive jail shell. \\
\command{exec <jail> [options] -- <argv>} & Execute a command with optional user,
working directory, and environment. \\
\command{run --name N --release R [--network NET] [--dns IP] -- <argv>} &
Run a command in an ephemeral jail, optionally overriding its repeatable
nameserver list; the jail is removed when the command exits. \\
\command{run --name N --release R [--network NET] [--dns IP] -d} &
Start an ephemeral jail and leave it running; takes no command. Use
\command{exec}, then \command{rm}. \\
\command{cp <source> <destination> [--preserve]} & Copy host-to-jail or
jail-to-host; use \command{jail:path} with a short or full jail name. \\
\command{rm <jails...> [--force] [--volumes]} & Remove ephemeral/managed jail
artifacts and optionally datasets; exits non-zero if any jail was not removed. \\
\command{logs <jail> [--follow] [--lines N]} & Read the jail log. \\
\bottomrule
\end{shiptable}

\section{Releases and builds}

\begin{shiptable}{L{0.39\textwidth}L{0.51\textwidth}}
\toprule
Command & Purpose \\
\midrule
\command{bootstrap <release> [--force] [--archives A] [--pkgbase]} & Acquire a
base userland. \\
\command{releases [list|delete|verify]} & Inspect or manage release roots. \\
\command{init [--file F] [--release R] [--toml] [--force]} & Generate a starter
Jailfile. \\
\command{build [--file F] [--name N] [--build-arg K=V] [--context D]
[--dry-run]} & Build an image root. \\
\command{template list|inspect|validate} & Work with Jailfiles and templates. \\
\command{commit <jail> <release>} & Freeze a jail into a reusable release. \\
\bottomrule
\end{shiptable}

\section{Networking and storage}

\begin{shiptable}{L{0.39\textwidth}L{0.51\textwidth}}
\toprule
Command & Purpose \\
\midrule
\command{network create <name> --subnet CIDR [options]} & Persist and realize a
named epair or netgraph network. \\
\command{network destroy <name> [--force]} & Remove a managed network. \\
\command{network list} & List managed network records and runtime state. \\
\command{expose <jail> --port P [options]} & Add a PF port forward. \\
\command{ports [jail]} & List persisted forwards. \\
\command{unexpose <jail>} & Remove a jail's forwards. \\
\command{snapshot create|list|rollback|delete} & Manage jail ZFS snapshots. \\
\command{clone <jail>@<snapshot> <name>} & Create a jail root from a snapshot. \\
\command{export <jail> [--output F] [--zfs-send]} & Produce a portable archive or
ZFS stream. \\
\command{import <file> [--name N] [--force] [--from manager]} & Restore native
or foreign jail data. \\
\command{migrate <jail> <user@host> [--remote-dataset D] [--keep]} & Send a
stopped ZFS-backed jail over SSH. \\
\bottomrule
\end{shiptable}

\section{Health and orchestration}

\begin{shiptable}{L{0.39\textwidth}L{0.51\textwidth}}
\toprule
Command & Purpose \\
\midrule
\command{health [jail] [--watch] [--interval N] [--json]} & Run configured
checks. \\
\command{stats [jail] [--json]} & Aggregate CPU, RSS, and process counts. \\
\command{supervise [--autostart]} & Run the warden for all configured jails;
the option reconciles and supervises only the effective boot set. \\
\command{qos <jail> --show|--startup|--steady} & Inspect or change a QoS phase. \\
\command{audit <jail> [--json]} & Render durable lifecycle history. \\
\command{armada init|up|down|build|ps|config} & Compose multi-jail workflows and
merged manifests. \\
\command{completion bash|elvish|fish|nu|zsh} & Generate shell completion. \\
\command{mcp [--listen ADDR:PORT]} & Serve MCP on stdio, or Streamable HTTP with \command{--listen} (token required off loopback). \\
\bottomrule
\end{shiptable}

\section{Machine-readable contract}

The global JSON format writes one versioned success object to standard output:

\begin{lstlisting}[language=JSON]
{
  "schema_version": 1,
  "ok": true,
  "data": []
}
\end{lstlisting}

Failures exit non-zero and write one object to standard error:

\begin{lstlisting}[language=JSON]
{
  "schema_version": 1,
  "ok": false,
  "error": {
    "code": "jail_not_found",
    "message": "Jail 'web' not found"
  }
}
\end{lstlisting}

Stable error codes are for program logic; messages are diagnostic. Commands
without a complete contract return \command{unsupported\_output\_format} rather
than mixing prose into JSON.

Read-only coverage includes status, validation, capabilities, the host planning
profile, networks, ports,
release listing, stats, one-shot health, audit, snapshot listing, and Armada
status. Lifecycle JSON is restricted to \command{--dry-run}. Streaming watch or
follow modes cannot fit a single response envelope.

\command{ps --host} returns the selected project and config path plus ownership
metadata for each row. It includes every manifest jail and running unmanaged
jails. Stopped unmanaged jails are not discoverable without a known manifest
and are omitted.

Legacy command-local \command{--json} flags remain unenveloped for compatibility.
New clients should prefer the global format and check
\setting{schema\_version} before decoding \setting{data}.

\section{Exit behavior}

Zero means the requested operation completed. Non-zero means parsing,
validation, escalation, host operation, or rendering failed. An in-jail
\command{exec} propagates the executed process status; treat that result as
workload data rather than a Blackship transport error in automation.
